% s-368-p368-m001; original PDF 368, 369
\Aphorism{s-368-p368-m001}{410}
Dem aristotelischen Satze zum Trotz, dass Existenz nie aus dem Begriff eines Dinges zu folgern ist (\Greek{τὸ} \Greek{εἶναι} \Greek{οὐκ} \Greek{οὐσία} \Greek{οὐδενί}), kann man das Dasein einer causal determinirten Welt auf diese Weise ontologisch ableiten, nämlich zeigen, dass ein vollkommen beliebiges Weltgeschehen sich von innen betrachtet als Causalverlauf darstellt. Es ist klar, dass die Vorstellung, einer Causalwelt anzugehören, eine augenblickliche Bewusstseinsillusion ist und als solche, bei chaotischem Zufallsspiel des transcendentalen Geschehens, nur vorübergehend und höchst selten auftauchen wird. Diese sporadischen Blitze erleuchten aber die ganze Strecke — der Intellect, der dieser Illusion fähig ist, wähnt sie \Sperr{beständig} erfüllt und aufrechterhalten zu finden. Warum? Dies eben liegt im Wesen der Illusion, im Begriff der determinirten Welt. Wenn wir stets in einem geordneten Kosmos zu leben glauben, so ist dies „stets“ ein Bestandtheil der fingirten Stetigkeit und Ordnung; eine Vorstellung wie die: „wir leben gegenwärtig in einer causalen Welt, werden aber mit nächstem aus ihr verstossen werden und dem Chaos anheimfallen“ — eine solche Vorstellung wäre antilogisch, also zum mindesten nicht die einer Causalwelt. Zu jeder Illusion gehören Nebenillusionen, die sie stützen, ergänzen, expliciren; zum Glauben an Causalität gehört der Glaube an Logik, und darum der Glaube an ewige, undurchbrochene, undurchbrechliche Causalität. Also, dass wir \Sperr{gelegentlich} an Causalität glauben, ist selbst in einer chaotischen Welt kein grösseres Wunder, als dass Einer unter Millionen Loosen das grosse Loos gewinnt; aus dem Begriff der Causalität aber, den wir in diesen bevorzugten Augenblicken entwerfen, folgt, dass wir \Sperr{ewig} in einer causalen Welt zu leben glauben. Dieses einfache Schema giebt die Deduction des Kosmos aus dem Chaos; dieses logische Netz ist das Sieb, das die chaotischen Zeitstrecken aus dem Bewusstsein herausfallen lässt und nur die kosmischen übrig behält.
\par

